
\section{Analog vs. Digital}
\label{section:analog_vs_digital}
\begin{frame}{Analog vs. Digital}%STARTCONTENT
    Bei der Informationsübertragung unterscheidet man grundsätzlich zwischen analogen und digitalen Verfahren.
    \begin{columns}
        \begin{column}{0.48\textwidth}
            \begin{itemize}
                \item<2-> \textbf{Analog}: kontinuierlich, beliebige Zwischenwerte
                \item<3-> \textbf{Digital}: in Stufen, nur bestimmte Werte, keine Werte dazwischen
            \end{itemize}
        \end{column}
        \begin{column}{0.48\textwidth}
            \only<2>{\begin{figure}
                \DARCimage{0.85\linewidth}{408include}
                \caption{\scriptsize Analoges Signal (kontinuierlich)}
                \label{n_digital_einleitung_analoges_signal}
            \end{figure}}%
            \only<3->{\begin{figure}
                \DARCimage{0.85\linewidth}{411include}
                \caption{\scriptsize Digitales Signal (abgestuft)}
                \label{n_digital_einleitung_digitales_signal}
            \end{figure}}%
        \end{column}
    \end{columns}
\end{frame}%ENDCONTENT